% !TeX root = ../papers/nc-p3.tex
\section{符号付きの六表示による全通常窓の八次判定}
\label{sec:rank-fifteen-oriented-eighth-cut}
\begin{theorem}
標数零のすべての通常な階数15窓において、四始点に共通な四次中心四本を選び、
八次まで中心単項式基底を取れる。各始点の八次で中心商は31次元、未切断経路は165次元である。
\end{theorem}
\begin{proof}
前節で残った非分岐単元対で、始点を保つために
定理\ref{thm:rank-fifteen-global-invariants}の座標 $(x,p,q)$ を使う。
$U=K^2$, $V=H^2$ と書くと
\[
 p=\frac{H^3K(U-9)}{(U-1)^3},\qquad p^2=x^3-27y,\qquad
 q^2=(1-x)^2-4y.
\]
剰余で $K=1/v$, $H=c(K-v)$, $X=c^2$ とすると $\bar p=c^3$ である。
そこで目標の $p$ に対して符号を含めて $c$ を固定する。$B=Y/X^3$ に対する
六つの $v$ のいずれでも、対応する $U$ と $K,H$ を既出のモニックな座標方程式で
持ち上げられる。この表示の $p'$ は $p'^2=p^2$ かつ $\bar p'=\bar p\ne0$ なので $p'=p$ である。
残る $q$ の符号は $K,H$ を保って選べる。実際 $n^2+(H/n)^2=U-1$ であり、
$n$ と $H/n$ の交換は $q$ の分子の因子 $2n^2+1-U$ の符号を反転する。
$q$ の他の因子はこの交換で変わらない。六根の符号は語の偶奇降下でどちらも含まれる。
従ってこの六表示のどれで得た始点0の証明書も、指定した窓へ\emph{添字を保って}
持ち上げられる。全十二表示を一緒に消去したときは、この最後の結論は直接には得られない。

第\ref{sec:rank-fifteen-bivariate-sixth-cut}節の第1,2,4,5列で始点0の右語を延ばすと、
八次の文字別次元は $(7,8,8,8)$ となる。
\path{computations/verify_rank_fifteen_oriented_eighth_cut.py} は原始多項式行に有理スカラー
$A/D$ を保持してこの消去を行う。原始行の非零係数 $b$ の有理主元としての分子・分母の次数和は
\[
 \deg A+\deg D+\deg b-2\deg\gcd(b,D)
\]
なので、第\ref{sec:rank-fifteen-bivariate-word-paths}節と同順位の主元を選べる。
除外多項式 $G(v,c)$ の次数は $(167,146)$ である。今度は $c$ を固定し、
$v^6=B+2v^4-v^2$ だけを使い、$G$ を $1,v,\ldots,v^5$ に関する六係数に分解する。
$B\ne0$ なので六表示は相異なり、全表示で除外される条件は六係数の同時消失である。
六係数の最大公約元は1である。$c,c-1,c+1,B,c^6B-(1-c^2)^2$ の単元因子を除き、
係数を低次側から二本ずつ組にして $B$ を消去する。三終結式から $c,c\pm1$ の因子を
除いた最大公約元の次数は33である。$c\in\mathbb F_9$ は既知なので、
より高次の既約因子の剰余体で六係数の $B$ に関する最大公約元を求める。
残る $(c,B)$ は六つの三次軌道と二つの四次軌道であり、個々の $B$ は
$\mathbb F_3(c)$ に属する。これらは元の四窓の有限体計算で四始点とも十三次まで確認する。
これでどの指定窓でも、少なくとも始点0において、八次商の上界31を実現できる。

最後に四中心の選択を共通化する。始点0での主張を一位置移動した窓にも適用する。
同じ $(x,y)$ 上の移動・反転は定理\ref{thm:rank-fifteen-rational-quotient}の符号変更で
表され、規格化テンソル上でも作用する。上で両方の $c$ と $q$ の符号を扱ったので、
添字を戻すと元の窓の各始点で少なくとも一つの四中心が取れる。
固定した窓の四次中心空間 $Z_4$ で、ある始点の八次商次元が31以下である条件は
$\operatorname{Gr}(4,Z_4)$ の開条件である。実際、普遍部分空間の局所基底を四式として、
関係行列の小行列式で書ける。各始点に対応する四つの非空開集合と、
前節の七次までの非空開集合はこの既約Grassmann多様体で同時に交わる。
標数零の基礎体上にも点がある。前節の下界31と生成写像の次元比較によって
八次の両次元と基底の主張が従う。
\end{proof}

\begin{corollary}
この四中心では十二次の商にも下界1がある。実際、中心単項式の生成により
$\dim Q_{p,12}\ge\binom{15}{3}-4\cdot31-10\cdot31-20\cdot1=1$ である。
従って九次から十二次までの目標上界 $20,10,4,1$ は、各次数で別々に証明してよい。
各上界の証明書が異なる四中心を用いる場合も、同じGrassmann多様体の開条件を交差させればよい。
\end{corollary}
